% GENERATED FILE. Edit canonical lessons, then run npm run generate:books.
% canonical-source-hash: fnv1a64:436d2fa99bfd8fd4
% canonical-lessons: TA-C432-nirnilai, TA-C432-kalvay, TA-C432-anaikkattu, TA-C432-kinaru, TA-C432-talumpu

\chapter{Canals and Wells}
\label{ch:ta-canals-and-wells}
\hlchaptermodality{432}

\begin{chapteropening}
\textbf{By the end of this chapter:} \emph{I can talk about bodies of water, canals, dams and wells, and water sloshing.}

Complete the last lesson of chapter 432: I can talk about bodies of water, canals, dams and wells, and water sloshing.
\end{chapteropening}

\section[\texorpdfstring{\ta{நீர்நிலை} (nīrnilai)}{நீர்நிலை (nīrnilai)}]{\ta{நீர்நிலை} (nīrnilai) — a body of water}

\label{lesson:TA-C432-nirnilai}



\begin{quote}
Before the new one: say the Tamil for to get scorched, then the Tamil for the water level.
\end{quote}

\subsection*{You'll want to know: \ta{நீர்நிலை}}
\textbf{\ta{நீர்நிலை}} — \emph{nīrnilai} — \textquotedblleft{}a body of water\textquotedblright{}.

\begin{quote}
\ta{ஏரிகளும்} \ta{குளங்களும்} \ta{முக்கியமான} \ta{நீர்நிலைகள்}.
\end{quote}

\begin{quote}
\emph{ērikaḷum kuḷaṅkaḷum mukkiyamāṉa nīrnilaikaḷ.}
\end{quote}

\begin{quote}
Lakes and ponds are important bodies of water.
\end{quote}

\ta{நீர்}, water, + \ta{நிலை}, a standing place.

\begin{cousinweb}
\textbf{\ta{நீர்நிலை}} (\emph{nīrnilai}) is built from \ta{நீர்}, \textquotedblleft{}water\textquotedblright{} (DEDR 3690), and \ta{நிலை}, \textquotedblleft{}a standing, a place\textquotedblright{} (DEDR 3675).
\end{cousinweb}

\begin{grammarlens}[title={What you've built}]
One more word for getting around.
\end{grammarlens}

\subsection*{Your turn}
\begin{itemize}
\raggedright
  \item \emph{Say it:} \emph{nīrnilai}
  \item \emph{Say it:} \emph{nīrnilai}, once more
  \item \emph{Say it:} \emph{ērikaḷum kuḷaṅkaḷum mukkiyamāṉa nīrnilaikaḷ.} — Lakes and ponds are important bodies of water.
  \item \emph{Recall:} say the Tamil for to get scorched, then the Tamil for the water level, then say \emph{nīrnilai} again
\end{itemize}

\begin{tcolorbox}[breakable,colback=teal!4,colframe=teal!35!black,title={Before you move on}]
What does \ta{நீர்நிலை} mean? (\textquotedblleft{}A body of water\textquotedblright{}.) Say it once more.
\end{tcolorbox}
\section[\texorpdfstring{\ta{கால்வாய்} (kālvāy)}{கால்வாய் (kālvāy)}]{\ta{கால்வாய்} (kālvāy) — a canal; a channel}

\label{lesson:TA-C432-kalvay}



\begin{quote}
Before the new one: say the Tamil for the water level, then the Tamil for a body of water.
\end{quote}

\subsection*{You'll want to know: \ta{கால்வாய்}}
\textbf{\ta{கால்வாய்}} — \emph{kālvāy} — \textquotedblleft{}a canal; a channel\textquotedblright{}.

\begin{quote}
\ta{கால்வாய்} \ta{வயல்களுக்குத்} \ta{தண்ணீர்} \ta{கொண்டு} \ta{போகிறது}.
\end{quote}

\begin{quote}
\emph{kālvāy vayalkaḷukkut taṇṇīr koṇṭu pōkiṟatu.}
\end{quote}

\begin{quote}
The canal carries water to the fields.
\end{quote}

The plural is \ta{கால்வாய்கள்}.

\begin{cousinweb}
\textbf{\ta{கால்வாய்}} (\emph{kālvāy}) is native Dravidian: DEDR 1480 lists it, \textquotedblleft{}an irrigation channel\textquotedblright{}.
\end{cousinweb}

\begin{grammarlens}[title={What you've built}]
One more word for getting around.
\end{grammarlens}

\subsection*{Your turn}
\begin{itemize}
\raggedright
  \item \emph{Say it:} \emph{kālvāy}
  \item \emph{Say it:} \emph{kālvāy}, once more
  \item \emph{Say it:} \emph{kālvāy vayalkaḷukkut taṇṇīr koṇṭu pōkiṟatu.} — The canal carries water to the fields.
  \item \emph{Recall:} say the Tamil for the water level, then the Tamil for a body of water, then say \emph{kālvāy} again
\end{itemize}

\begin{tcolorbox}[breakable,colback=teal!4,colframe=teal!35!black,title={Before you move on}]
What does \ta{கால்வாய்} mean? (\textquotedblleft{}A canal; a channel\textquotedblright{}.) Say it once more.
\end{tcolorbox}
\section[\texorpdfstring{\ta{அணைக்கட்டு} (aṇaikkaṭṭu)}{அணைக்கட்டு (aṇaikkaṭṭu)}]{\ta{அணைக்கட்டு} (aṇaikkaṭṭu) — a dam}

\label{lesson:TA-C432-anaikkattu}



\begin{quote}
Before the new one: say the Tamil for a body of water, then the Tamil for a canal.
\end{quote}

\subsection*{You'll want to know: \ta{அணைக்கட்டு}}
\textbf{\ta{அணைக்கட்டு}} — \emph{aṇaikkaṭṭu} — \textquotedblleft{}a dam\textquotedblright{}.

\begin{quote}
\ta{ஆற்றின்} \ta{குறுக்கே} \ta{ஒரு} \ta{பெரிய} \ta{அணைக்கட்டு} \ta{இருக்கிறது}.
\end{quote}

\begin{quote}
\emph{āṟṟiṉ kuṟukkē oru periya aṇaikkaṭṭu irukkiṟatu.}
\end{quote}

\begin{quote}
There is a big dam across the river.
\end{quote}

\ta{அணை}, an embankment, + \ta{கட்டு}, to build.

\begin{cousinweb}
\textbf{\ta{அணைக்கட்டு}} (\emph{aṇaikkaṭṭu}) is built from \ta{அணை}, \textquotedblleft{}an embankment, a dam\textquotedblright{} (DEDR 122), and \ta{கட்டு}, \textquotedblleft{}to tie, to build\textquotedblright{} (DEDR 1147).
\end{cousinweb}

\begin{grammarlens}[title={What you've built}]
One more word for getting around.
\end{grammarlens}

\subsection*{Your turn}
\begin{itemize}
\raggedright
  \item \emph{Say it:} \emph{aṇaikkaṭṭu}
  \item \emph{Say it:} \emph{aṇaikkaṭṭu}, once more
  \item \emph{Say it:} \emph{āṟṟiṉ kuṟukkē oru periya aṇaikkaṭṭu irukkiṟatu.} — There is a big dam across the river.
  \item \emph{Recall:} say the Tamil for a body of water, then the Tamil for a canal, then say \emph{aṇaikkaṭṭu} again
\end{itemize}

\begin{tcolorbox}[breakable,colback=teal!4,colframe=teal!35!black,title={Before you move on}]
What does \ta{அணைக்கட்டு} mean? (\textquotedblleft{}A dam\textquotedblright{}.) Say it once more.
\end{tcolorbox}
\section[\texorpdfstring{\ta{கிணறு} (kiṇaṟu)}{கிணறு (kiṇaṟu)}]{\ta{கிணறு} (kiṇaṟu) — a water well; a well}

\label{lesson:TA-C432-kinaru}



\begin{quote}
Before the new one: say the Tamil for a canal, then the Tamil for a dam.
\end{quote}

\subsection*{You'll want to know: \ta{கிணறு}}
\textbf{\ta{கிணறு}} — \emph{kiṇaṟu} — \textquotedblleft{}a water well; a well\textquotedblright{}.

\begin{quote}
\ta{கிராமத்தில்} \ta{ஒரு} \ta{பழைய} \ta{கிணறு} \ta{இருக்கிறது}.
\end{quote}

\begin{quote}
\emph{kirāmattil oru paḻaiya kiṇaṟu irukkiṟatu.}
\end{quote}

\begin{quote}
There is an old well in the village.
\end{quote}

Plural: \ta{கிணறுகள்}, \emph{kiṇaṟukaḷ}.

Before a case the \ta{று} doubles: \ta{கிணற்றில்}, in the well.

\begin{cousinweb}
\textbf{\ta{கிணறு}} (\emph{kiṇaṟu}) is native Dravidian (DEDR 1998), \textquotedblleft{}a well\textquotedblright{}.
\end{cousinweb}

\begin{grammarlens}[title={What you've built}]
One more word for getting around.
\end{grammarlens}

\subsection*{Your turn}
\begin{itemize}
\raggedright
  \item \emph{Say it:} \emph{kiṇaṟu}
  \item \emph{Say it:} \emph{kiṇaṟu}, once more
  \item \emph{Say it:} \emph{kirāmattil oru paḻaiya kiṇaṟu irukkiṟatu.} — There is an old well in the village.
  \item \emph{Recall:} say the Tamil for a canal, then the Tamil for a dam, then say \emph{kiṇaṟu} again
\end{itemize}

\begin{tcolorbox}[breakable,colback=teal!4,colframe=teal!35!black,title={Before you move on}]
What does \ta{கிணறு} mean? (\textquotedblleft{}A water well; a well\textquotedblright{}.) Say it once more.
\end{tcolorbox}
\section[\texorpdfstring{\ta{தளும்பு} (taḷumpu)}{தளும்பு (taḷumpu)}]{\ta{தளும்பு} (taḷumpu) — to slosh; to brim over}

\label{lesson:TA-C432-talumpu}



\begin{quote}
Before the new one: say the Tamil for a dam, then the Tamil for a water well.
\end{quote}

\subsection*{You'll want to know: \ta{தளும்பு}}
\textbf{\ta{தளும்பு}} — \emph{taḷumpu} — \textquotedblleft{}to slosh; to brim over\textquotedblright{}.

\begin{quote}
\ta{வண்டி} \ta{குலுங்கியதால்} \ta{வாளியில்} \ta{தண்ணீர்} \ta{தளும்பியது}.
\end{quote}

\begin{quote}
\emph{vaṇṭi kuluṅkiyatāl vāḷiyil taṇṇīr taḷumpiyatu.}
\end{quote}

\begin{quote}
Because the cart jolted, the water sloshed in the bucket.
\end{quote}

Past (it …): \ta{தளும்பியது}, \emph{taḷumpiyatu}.

Of eyes: \ta{கண்ணீர்} \ta{தளும்பியது}, eyes brimmed with tears.

\begin{cousinweb}
\textbf{\ta{தளும்பு}} (\emph{taḷumpu}), \textquotedblleft{}to slosh, to brim over\textquotedblright{}, is a Tamil verb whose origin is not settled.
\end{cousinweb}

\begin{grammarlens}[title={What you've built}]
One more word for getting around.
\end{grammarlens}

\subsection*{Your turn}
\begin{itemize}
\raggedright
  \item \emph{Say it:} \emph{taḷumpu}
  \item \emph{Say it:} \emph{taḷumpu}, once more
  \item \emph{Say it:} \emph{vaṇṭi kuluṅkiyatāl vāḷiyil taṇṇīr taḷumpiyatu.} — Because the cart jolted, the water sloshed in the bucket.
  \item \emph{Recall:} say the Tamil for a dam, then the Tamil for a water well, then say \emph{taḷumpu} again
\end{itemize}

\begin{tcolorbox}[breakable,colback=teal!4,colframe=teal!35!black,title={Before you move on}]
What does \ta{தளும்பு} mean? (\textquotedblleft{}To slosh; to brim over\textquotedblright{}.) Say it once more.
\end{tcolorbox}
